% Neutral A4 single-column preprint adaptation; compile with XeLaTeX.
% All prose, authors, affiliations, references, code and the figure are original placeholders.
% No real source-paper content, results, figures, identifiers, dates or logos.
\PassOptionsToPackage{no-math}{fontspec}
\documentclass[UTF8,fontset=none,scheme=plain]{ctexart}
\usepackage[paperwidth=595.276bp,paperheight=841.89bp,left=62.36bp,right=60.876bp,top=49bp,bottom=48bp,headheight=12bp,headsep=9bp,footskip=22bp]{geometry}
% Shared dependencies and neutral material. Compile with XeLaTeX.
% No publisher logos, paper text, data, figures, or download metadata are included.
\usepackage{amsmath,amssymb,graphicx,xcolor,tikz,fancyhdr,needspace,xurl,hyperref,balance,etoolbox}
% The legacy balance package omits the first split box's final baseline depth.
% Reserve that measured depth rather than hiding a CJK overfull-box warning.
\makeatletter
\patchcmd{\@BAlancecol}{\setbox\@leftcolumn=\vbox to \@tempdima}{\advance\@tempdima by \dp1\relax\setbox\@leftcolumn=\vbox to \@tempdima}{}{\PackageWarning{cns-template}{Balance depth patch unavailable; inspect final columns}}
\makeatother
\usetikzlibrary{arrows.meta,positioning}
\usepackage[protrusion=true,expansion=false]{microtype}
\IfFontExistsTF{Noto Serif CJK SC}{%
  \setCJKmainfont{Noto Serif CJK SC}%
}{%
  \setCJKmainfont[BoldFont=FandolSong-Bold.otf,ItalicFont=FandolKai-Regular.otf]{FandolSong-Regular.otf}%
}
\IfFontExistsTF{Noto Sans CJK SC}{%
  \setCJKsansfont{Noto Sans CJK SC}%
}{%
  \setCJKsansfont[BoldFont=FandolHei-Bold.otf]{FandolHei-Regular.otf}%
}
\setmainfont{TeX Gyre Termes}\setsansfont{TeX Gyre Heros}
\setmonofont{TeX Gyre Cursor}
\definecolor{JournalBlue}{HTML}{0097C9}
\definecolor{JournalRed}{HTML}{B83443}
\definecolor{ScienceRed}{HTML}{CD1021}
\definecolor{CiteBlue}{HTML}{0000EE}
\hypersetup{colorlinks=true,linkcolor=JournalBlue,urlcolor=JournalBlue,
  pdftitle={CNS论文汉化：中性版式示例},pdfauthor={示例模板；非期刊官方模板},
  pdfsubject={Neutral source-inspired Chinese layout demonstration; no real research claims}}
\urlstyle{same}\raggedbottom
\setlength{\parindent}{1em}\setlength{\parskip}{0pt}
\setlength{\emergencystretch}{2em}\tolerance=2200
\XeTeXlinebreaklocale "zh"\XeTeXlinebreakskip=0pt plus .1pt
\pagestyle{fancy}\fancyhf{}
\renewcommand{\headrulewidth}{0pt}\renewcommand{\footrulewidth}{0pt}
\newcommand{\PaperTitle}{中文题名占位：从组织结构到证据解释}
\newcommand{\SampleAuthors}{Author A\textsuperscript{1}, Author B\textsuperscript{2,*}}
\newcommand{\SampleAffiliations}{\textsuperscript{1}作者单位一（示例占位）。\quad
  \textsuperscript{2}作者单位二（示例占位）。\par
  *通信作者信息待填。DOI、收稿和出版日期均应从目标原文核实。}
\newcommand{\SampleAbstract}{本段是中性摘要占位文字，不陈述任何真实研究发现。它演示中文题名、作者信息、摘要和双栏正文之间的层级关系。实际使用时，应完整替换为目标论文的忠实中文译文，并保留研究对象、方法、结果和结论的限定条件。图注、参考文献、作者声明和来源信息均属于独立的覆盖项。模板中的图形仅说明排版结构，不表示生物学机制、实验结果或统计关系。}
\newcommand{\DemoCaption}{\textbf{图1. 中性证据链示意（占位图，无实验数据）。}
  \textbf{A}，输入节点表示需要核对的原文材料。
  \textbf{B}，处理节点表示逐段转换及复核。
  \textbf{C}，输出节点表示正文、图注和书目等独立交付部分。
  箭头仅表示本示例的阅读顺序，不代表因果关系、效应方向或数量变化。
  该图由模板自行绘制，所有节点均为文字占位；无样本量、误差线、统计检验或比例尺。
  实际论文图注应逐项保留原稿报告的样本量、单位、检验方法、限定及图形署名。}
\newcommand{\DemoVector}{%
 \resizebox{\linewidth}{!}{%
 \begin{tikzpicture}[font=\sffamily\small,node distance=18mm,
   box/.style={draw=black!70,rounded corners=2pt,fill=black!3,minimum width=26mm,minimum height=16mm,align=center},
   flow/.style={-{Stealth[length=2.5mm]},line width=.7pt}]
  \node[box] (a) {\textbf{A}\quad 输入\\原文材料};
  \node[box,right=of a] (b) {\textbf{B}\quad 处理\\转换与复核};
  \node[box,right=of b] (c) {\textbf{C}\quad 输出\\中文译稿};
  \draw[flow] (a)--(b);\draw[flow] (b)--(c);
  \node[below=11mm of a,align=center] {标题与摘要\\正文与图表};
  \node[below=11mm of b,align=center] {保留限定\\检查单位};
  \node[below=11mm of c,align=center] {可追溯引文\\明确原文信息};
  \node[below=29mm of b,font=\sffamily\footnotesize,text=black!65] {示意占位；不包含真实实验数据};
 \end{tikzpicture}}%
}
\newcommand{\SampleBibliography}{%
 \section*{参考文献（示例占位）}
 {\fontsize{7.6bp}{10bp}\selectfont\setlength{\parindent}{0pt}\interlinepenalty=10000
 \hypertarget{ref:1}{}1. Author A. (20XX). Example reference title. \textit{Journal placeholder}. 卷、页码与DOI待填。此条仅供演示引文跳转。\par\smallskip
 \hypertarget{ref:2}{}2. Author B. (20YY). Another example reference title. \textit{Journal placeholder}. 出版信息待填；不要把占位条目带入正式译稿。\par}
}
\newcommand{\SampleSourceInfo}{%
 \section*{原文信息}
 以下字段全部为示例占位：作者；英文原标题；期刊；年份；卷期；页码或文章号；DOI或来源链接。实际使用时必须逐项替换并核对。\par
 本文件是\JournalName 风格的非官方中文译稿样板，依据一次已核验文章版式进行结构化改编，不代表该期刊所有年代或文体的统一标准，也不保证像素级复刻。无真实论文内容或期刊标识图像。\par
}
% Deliberately neutral demo prose, not translated from a published paper.
\newcommand{\SampleOpening}{%
 \section*{引言：示例正文与引文}
 这是一段用于检查双栏阅读节奏的中性示例文字。段落中的术语、缩写和标点需要在实际项目中与源文保持一致。中文译稿应呈现原论文的完整论述，而不是把多段证据压缩成一段摘要。这里的引文只链接到本文件末尾明确标为占位的参考文献\SampleCiteOne。\par
 若原句仅表示相关、可能或提示，译文也应保留相同的证据强度。数字、单位、正负号、上下标和范围端点应逐项核对。原文存在内部矛盾时，应分别保留相应位置的陈述，再用简短译注说明；不能根据常识把不同的数字静默统一。\par
 \subsection*{内容层级与跨页衔接}
 此小标题演示二级层级。实际论文可能采用不同的小标题密度，模板不会替代对原文结构的识别。双栏文章需按左栏、右栏和跨页续句的顺序读取。若段落在页末中断，应在完成整句核对后接入下一页，避免把图中文字误当正文，或重复翻译同一段。\par
 图表具有独立的信息范围。正文引用到某幅图，并不意味着其图注可以省略。图注中的每个分面、比较对象、时间点、样本量及统计说明都应对应源文。英文图内标签可以保留，并在图旁或文末提供中文释义；实际采用哪种安排，需要依据图形尺寸和译文长度判断\SampleCiteTwo。\par
 \subsection*{来源与可维护数据}
 本段继续演示正文密度。可编辑文本和图注数据应与最终版面分开维护，以便修订时准确定位变化。若只调整版式，应能证明正文、图注和书目内容没有被意外改写。更换字体或纸张后，须重新检查断行、图注与图片的相对位置以及页脚可读性。\par
 文内引用与书目之间应有明确对应关系。这里的两条演示书目没有真实DOI，也没有虚构的实验信息。使用模板时可增删条目，但必须同时更新编号、链接目标和引用位置。页码只表示中文译稿自己的分页，不用于冒充原出版页码。\par
 \subsection*{结果叙述的演示占位}
 本节的文字仅用于呈现长段落在双栏页面中的流动，不报告真实结果。实际翻译应保留观察对象、实验条件和结局指标之间的关系。例如，原文分别描述体外实验与动物实验时，不能把两类结果合并为同一层面的结论；研究使用的对照、时间顺序和分析范围也必须保留。图表中的比较范围不应扩大为对整个人群或全部疾病阶段的判断。\par
 论述中的转折与限制同样属于正文。若研究先提出一种解释，再讨论另一种可能，译文应沿用这一推进顺序。作者明确写出的不确定性应能在中文中被识别；如果尚未确定某个现象的方向或机制，就不能改成已经建立的因果关系。对于术语近似但研究层级不同的概念，可在首次出现时保留缩写，后文保持统一。\par
 \subsection*{方法信息与可复核性}
 方法部分往往包含密集数字和专有名词。翻译时须区分浓度、总量、频率和处理时长，不能因为中文句式调整而改变参数的归属。试剂名称、软件版本和数据访问标识通常保留原样；上下限、排除规则和显著性标准则需要逐字符核对。这里没有提供任何真实实验参数，避免把模板占位误认为可以直接执行的方法。\par
 表格与补充材料的引用也需要保留原层级。正文对补充图的指向并不意味着模板中已经附有该补充图；若源文件只包含主文，应如实标注所交付的范围。图表编号、附录顺序和参考文献数量都应由目标文章决定，本示例中的结构不构成固定配额，也不要求所有论文采用同样的章节。\par
}
\newcommand{\SampleClosing}{%
 \section*{讨论：版式适配与范围}
 上方矢量图是模板原创的排版占位图，不来自任何论文。图形与完整演示图注被放在同一页，便于读者直接核对分面说明。真实项目应使用经过授权或用户提供的原图素材，并保留图注署名与许可信息。尺寸不足时，可适度缩放或调整栏宽，不以删除图注解决空间问题。\par
 模板仅记录一次已核验文章的可观察设计特征。即使文章出自同一期刊，不同年代、研究论文与综述、校正清样与正式版本也可能采用不同的页眉和引文形式。应以用户提供的实际PDF为准，不能仅凭期刊名称推断配色、字号或引用方式\SampleCiteOne。\par
 \subsection*{演示检查项}
 本样板用于展示可编译性、中文字形、标题层级、跨栏图注、页眉页脚和内部链接。图中没有实际数据，因此不应被当作科学插图引用。替换内容后，还需检查所有页的渲染结果，确认无截字、碰撞、异常空白和遗漏段落。\par
 \SampleBibliography
 \SampleSourceInfo
}


\usepackage{fvextra,booktabs}
\renewcommand{\familydefault}{\sfdefault}
\renewcommand{\CJKfamilydefault}{\CJKsfdefault}
\setCJKmonofont{FandolFang-Regular.otf}
\renewcommand{\normalsize}{\fontsize{9.5bp}{13.5bp}\selectfont}
\setlength{\parindent}{0pt}\setlength{\parskip}{5bp}
\definecolor{PreprintBlue}{HTML}{0000FF}
\hypersetup{colorlinks=true,linkcolor=PreprintBlue,urlcolor=PreprintBlue,
  pdftitle={A4单栏预印本中性中文样板},pdfauthor={原创中性版式演示},
  pdfsubject={Neutral single-column preprint demonstration; no source research content}}
\ctexset{section={format=\sffamily\bfseries\fontsize{12bp}{16bp}\selectfont,beforeskip=15bp,afterskip=6bp},
  subsection={format=\sffamily\bfseries\fontsize{10.3bp}{14bp}\selectfont,beforeskip=11bp,afterskip=4bp}}
\fancyhead{}\renewcommand{\headrulewidth}{0pt}
\renewcommand{\footrulewidth}{.35bp}
\fancyfoot[L]{\sffamily\fontsize{7bp}{9bp}\selectfont 预印本样板｜非官方中文译稿}
\fancyfoot[R]{\sffamily\fontsize{7bp}{9bp}\selectfont 中文译页\ \thepage}
\fvset{fontsize=\footnotesize,breaklines=true,breakanywhere=true,frame=single,rulecolor=\color{black!15},framesep=5bp}
\newcommand{\PreprintCite}[1]{\textsuperscript{\hyperlink{preprint-ref:#1}{#1}}}
\newcommand{\PreprintCaption}[1]{\par\begingroup\fontsize{8.5bp}{11.5bp}\selectfont #1\par\endgroup}
\begin{document}\normalsize
{\bfseries\fontsize{20bp}{25bp}\selectfont 中文题名占位：完整内容的单栏呈现与复核\par}
\vspace{4bp}
示例作者甲\textsuperscript{1,*}、示例作者乙\textsuperscript{1}、示例作者丙\textsuperscript{2,*}\par
{\fontsize{8.5bp}{11.5bp}\selectfont
\textsuperscript{1}示例单位一（原创占位）。\quad\textsuperscript{2}示例单位二（原创占位）。\par
*贡献及通信标记应按目标原文填写；本样板不提供真实邮箱。\par}
\vspace{13bp}
{\bfseries
本段为连续摘要的原创中性占位，不报告真实研究结果。它展示预印本首页全宽题名、作者单位、加粗摘要和同页单栏正文之间的层级。完整译稿应忠实保留研究问题、方法、结果和限定，而不是用内容概括替代逐段翻译。预印本的服务器、版本、发布日期、标识符及许可，应分别以实际源文件所载内容为依据；文件属性中的生成时间不等同于公开发布日期。本样板不借用任何期刊标识，也不为缺失的 DOI、卷期或正式页码填入推测值。正文与补充材料的范围以用户提供文件为准，中文可依长度自然重新分页。\par}
\vspace{9bp}
\section*{引言}
本页演示黑白无衬线单栏正文及蓝色数字上标\PreprintCite{1}。单栏预印本可能具有比期刊论文更长的方法和补充材料；内容清单应涵盖正文、完整图注、公式、伪代码、声明及各套参考文献。正文中的跨页续句需要先恢复完整，再转换为中文。原文没有节标题的段落不应为方便排版而任意改写其逻辑结构。

可编辑内容与排版代码应分别维护。调整字号或图像位置时，应保留正文、数值、引用和公式的内容指纹；对已核对译文换版式，不应重新翻译造成内容漂移。本段仅用于展示连续段落之间的正常留白及首行不缩进，不是某一研究的背景或结论。

\subsection*{范围与版本}
预印本身份、正式出版状态和同行评审状态是不同信息。译稿应明确自己所依据的版本，不把后来发表版本的信息无说明地混入当前来源。若原文件未给出服务器或 DOI，应说明未载明，而不根据标题相似性补造。此处所有版本字段均为中性占位，不对应任何真实预印本。

对图中保留的英文，完整中文释义应覆盖面板标题、坐标轴、图例及整句说明。基因名、模型名、数据标识和科学符号可按原样保留，但不能用零散术语表代替应译的说明句。每幅图的边界应单独核对，图注允许在必要时自然跨页\PreprintCite{2}。

\subsection*{完整性记录}
内容覆盖以源段和内容对象为单位，而非仅比较页数。正文、方法、补充说明和参考文献可采用各自稳定的标识；主文与补充书目重置编号时，还需要独立链接命名空间，避免同号误链。本页的引文仅指向末页两个明确标注的示例条目。

\clearpage
\section*{方法与图文示例}
本页演示可续页的代码、独立编号公式及原创矢量图。真实伪代码中的变量、数组维度、运算符和缩进应逐项保留；公式需对照源页检查上下标、求和范围与分母。下面的代码和公式仅说明排版方法，不是从任何论文提取的算法。
\subsection*{伪代码}
\begin{Verbatim}
def prepare_document(blocks, coverage):
    output = []
    for block in blocks:
        assert block.source_id in coverage
        output.append(block.verified_translation)
    return output
\end{Verbatim}
代码使用可换行的等宽字体环境，真实长代码块允许跨页，不整体包入不可断页盒。源文件中的注释若表达方法含义，也需翻译或给出完整中文对应。本例变量仅代表文档对象，不含真实实验参数。

\subsection*{公式与符号}
若 $a_i$ 是第 $i$ 个演示段落的字符数，$n$ 为演示段落数，则算术平均字符数写为
\begin{equation}
\bar a=\frac{1}{n}\sum_{i=1}^{n}a_i.\tag{1}
\end{equation}
该式为原创基础排版示例，不报告数值或科研发现。真实公式编号应跟随原文；若源文有疑点，须与译者修正清楚区分。

\vspace{8bp}
\begin{center}\begin{minipage}{.94\textwidth}\DemoVector\end{minipage}\end{center}
\PreprintCaption{\textbf{图1｜内容处理顺序（原创中性占位）。}
\textbf{A}，输入节点表示题名、摘要、正文、图表及附录等源内容。
\textbf{B}，处理节点表示恢复阅读顺序、逐段转换及核对限定条件。
\textbf{C}，输出节点表示完整中文译稿、可编辑内容和覆盖记录。
箭头仅指本示例的阅读顺序，不表示生物学因果关系、实验效应、统计差异或数值变化。图由模板自行绘制，无真实数据、原图或官方标识。图注与正文采用可正常断页的段落；真实长图应按其面板、长宽及图注长度单独安排。}

\subsection*{版面检查}
正文版心与宽图版心应分别测量。图可以按目标源稿适当伸入正文两侧留白，但不得截断坐标轴、面板字母或图例。图像应优先保留干净矢量素材；下载来源信息需实际移除，而非仅遮白或更改裁切框。

\clearpage
\section*{补充材料与覆盖示例}
补充图、补充说明、补充表描述和补充参考文献应按提供文件独立清点。若源 PDF 只列出表格名称与说明，未包含所指外部表格数据，应翻译已经提供的描述，并如实注明数据表未随该文件给出。本段为内容范围示例，不表示本样板附带任何实验附件。

\subsection*{补充说明（原创中性占位）}
跨页续段只能合并一次。图中英文保留时，应在译稿中以按图或面板分组的方式提供中文标签对应；图注内容与标签释义分别保留，不因相似而省略。对包含多个独立参考文献列表的文件，应记录各列表的编号起点和引用对象，并逐条检查内部跳转。

\subsection*{覆盖记录（原创中性占位）}
\begin{tabular}{@{}p{.19\textwidth}p{.43\textwidth}p{.30\textwidth}@{}}
\toprule
内容对象 & 需要保留的边界 & 演示校核状态\\\midrule
正文 & 段落、标题、引用、跨页续句 & 按源内容逐项记录\\[4bp]
图与图注 & 全部面板及完整说明 & 单独核对裁区与文字\\[4bp]
方法 & 参数、公式、伪代码和限定 & 对照源页逐项复核\\[4bp]
补充材料 & 实际提供的附图、说明及表格 & 区分已提供与未提供\\[4bp]
\bottomrule
\end{tabular}
\par\smallskip
以上字段为原创文档检查示例，不是科学数据、统计结果或期刊报告标准。

\section*{参考文献（中性占位）}
{\fontsize{8.3bp}{11.3bp}\selectfont
\hypertarget{preprint-ref:1}{}1. Example author. Example document title. Source and version details to be replaced from the supplied file.（原创占位，非真实文献。）\par
\hypertarget{preprint-ref:2}{}2. Example author. Another example source title. Publication and identifier fields to be verified.（原创占位，非真实文献。）\par}

\section*{声明（中性占位）}
致谢、资助、作者贡献、利益冲突、数据可用性、代码可用性及许可均应按实际原稿逐项翻译。本文件不包含真实研究、资助、数据或作者贡献；此处只演示声明区的层级与位置。

\section*{原文信息}
作者、英文题名、预印本服务器、版本、明确载明的发布日期、DOI 或其他标识符、页数及许可信息，均须从目标来源核实并填写。原稿未列字段应说明未载明。预印本通常不具有期刊卷期；不能将文件生成日期当作发布日期。

本文件为 A4 单栏预印本的非官方中性中文样板，不属于某期刊官方模板。全部文字、示意图、代码、表格、作者单位和书目均为原创占位，三页人工分页仅用于演示，不限制完整译稿长度。仅供学术交流、出版物版权归原作者及出版商所有。
\end{document}
